\documentclass[12pt,a4paper]{article}

% -----------------------------------------------------------------------
% Packages
% -----------------------------------------------------------------------
\usepackage[utf8]{inputenc}
\usepackage[T1]{fontenc}
\usepackage{lmodern}
\usepackage[english]{babel}
\usepackage{amsmath,amssymb,amsfonts}
\usepackage{booktabs}
\usepackage{graphicx}
\usepackage{hyperref}
\usepackage{cleveref}
\usepackage{xcolor}
\usepackage{listings}
\usepackage{algorithm}
\usepackage{algpseudocode}
\usepackage{tikz}
\usepackage{pgfplots}
\pgfplotsset{compat=1.18}
\usepackage{subcaption}
\usepackage{float}
\usepackage{multirow}
\usepackage{enumitem}
\usepackage{natbib}
\bibliographystyle{plainnat}

% Code listings style
\lstset{
  basicstyle=\ttfamily\small,
  breaklines=true,
  frame=single,
  numbers=left,
  numberstyle=\tiny\color{gray},
  keywordstyle=\color{blue},
  commentstyle=\color{green!60!black},
  stringstyle=\color{red!70!black},
}

\hypersetup{
  colorlinks=true,
  linkcolor=blue!70!black,
  citecolor=green!50!black,
  urlcolor=blue!60!black,
}

% -----------------------------------------------------------------------
% Title
% -----------------------------------------------------------------------
\title{\textbf{Shut Your Eyes and C}\\[0.3em]
\large Programming Languages, Jungian Archetypes,\\
and the Architecture of Machine Understanding}

\author{Victor Fermino}


\date{\today}

\begin{document}
\maketitle

\begin{abstract}
We present \textsc{DisattentionFormer}, a Transformer architecture that
embeds Jungian archetypal tensors---pre-computed covariance structures
derived from mythic and literary passages---as geometric priors that
deform the attention space. Trained on an intentionally ordered corpus
of 19~texts spanning from Genesis to Joyce's \textit{Ulysses}, the model
treats corpus order as architectural structure (\texttt{shuffle=False}).
We compare \textsc{DisattentionFormer} ($\approx$27.7M parameters) against
a parameter-matched vanilla Transformer baseline on perplexity, generation
quality, and archetype activation. Despite its rich philosophical
scaffolding, the model operates by the same mechanism as every other
model in the compilatory chain: faithful translation without comprehension.
It activates archetypal tensors via lexical proximity, not conceptual
recognition. It cannot see that the Shadow is the Persona inverted,
because that requires understanding, and parsers do not understand. The
work situates itself at the intersection of programming language theory
(Knuth's Literate Programming), analytical psychology (Jung's archetypal
theory), and deep learning, arguing that neural architectures---like C
compilers, like obfuscated code---are grammars that constrain expression
without generating meaning, and that \emph{understanding this} is more
valuable than prompting a black box to produce plausible text.
\end{abstract}

\section{Introduction: The Act of Translation}
\label{sec:introduction}

% NOTE: This section incorporates and expands the content from the
% original "Shut Your Eyes and C" manuscript (pages 1--3).
% The user's existing text about Hello World in C, obfuscated C,
% and the compiler-as-translator metaphor should be placed here,
% followed by the transition into ML architectures.

The act of developing software is an act of translation. When someone
writes a program that says \texttt{"Hello world"}, what they are doing
is finding an interface to meet the machine at the halfway point---just
like when one learns how to flex their tongue and extend a vowel to say
\textit{Ol\'a} in Portuguese as opposed to \textit{Hola} in Spanish.

\begin{lstlisting}[language=C, caption={The canonical greeting.}]
#include <stdio.h>
int main() {
    printf("Hello world");
    return 0;
}
\end{lstlisting}

One might wonder what is happening. Is it bureaucracy from a bygone era
of basement-dwelling programmers who wanted everything to be complicated?
If we take this apparent verbosity in context, however, we can clearly see
the attempt to make everything clear---a way of making things
\emph{communicable} to a machine, or to a compiler, which is just a
translator.

\texttt{\#include} instructs the compiler to include a different library;
\texttt{<stdio.h>} is the standard input/output library, so the compiler
knows that this program depends on a terminal window. The \texttt{main()}
method is the main thing that will happen. The \texttt{int} return
condition means the method expects an integer at the end: zero means
success; any other number signals an error.

We could say that this whole process of displaying a few words on a black
screen is too complex, but the structure itself is somewhat similar to how
we say \textit{Ciao!}\ in Italian: there is a social context for when we
are going to say something (the inclusion of a library for standard
communication), then the act of speech (a method), and then the utterance
itself.

\subsection{Obfuscation as Poetics}

There are competitions such as the International Obfuscated C Code Contest
(IOCCC, since 1984) that award code which is \emph{creatively} obfuscated.
Consider a program that calculates $\pi$ by looking at its own area---a
circle of underscores whose macro definition hides the counting logic. The
code is simultaneously valid C \emph{and} a visual poem. This is not
unlike how Joyce's \textit{Finnegans Wake} is simultaneously valid English
\emph{and} something else entirely: a portmanteau language that demands the
reader become a compiler.

\subsection{From Compilation to Generation}

All of this might seem too deterministic, and it is, when it comes to basic
translations. Any LLM, when prompted with a ``Hello world,'' will respond
differently. But let us consider \emph{temperature}, or better yet,
\emph{entropy}, which is what makes current AI paradigms what they are known
for.

A Transformer is, at its core, a text prediction system. If we write
``The author of \textit{Ulysses} is James \_\_\_\_\_,'' a well-trained
model will have ``Joyce'' as the most likely outcome. But it will also
predict the next few tokens---usually an explanation or a biography.
Current LLM architectures have multiple steps to turn mere prediction
into something that \emph{feels} alive: they usually do not pick the most
likely outcome, and deciding the weight is basically prioritising some
predictions over others.

This raises a question that is as much linguistic as it is computational:
\textbf{what if the architecture itself could carry meaning?} What if,
instead of treating the Transformer as a neutral vessel for statistical
patterns, we designed its components to embody a specific theory of how
meaning is structured?

And, more provocatively: what would it reveal if the answer were
\emph{not much}?

This paper pursues both questions simultaneously. We build a Transformer
whose architecture encodes Jungian archetypal theory as geometric
structure. We give it every advantage: hand-crafted curvature tensors,
an intentionally ordered mythic corpus, auxiliary losses designed to
maximise internal differentiation. And we show that, despite all this,
the model still operates at the level of lexical fingerprints---activating
``Great Father'' for Genesis not because it grasps the concept of
patriarchal divinity but because the words overlap with its seed passages.
It cannot recognise that the Shadow is the Persona inverted, because
\emph{that requires understanding}, and parsers do not understand.

The argument is not that the architecture fails. The argument is that its
\emph{mode of success}---statistical pattern matching within a structured
geometric bias---is the same mode of success at every level of the
compilatory chain, from \texttt{gcc} translating \texttt{printf} to a
syscall, to GPT-4 translating a prompt into plausible prose. The
temperature, the nucleus sampling, the apparent creativity---it is
obfuscated C all the way down.

% ======================================================================
\section{Background}
\label{sec:background}

\subsection{Programming Languages as Communication}
\label{sec:literate}

Knuth's Literate Programming \citep{knuth1984literate} proposed that
programs should be written primarily for human understanding, with machine
execution as a secondary concern. The program becomes a \emph{text}, and
the act of programming becomes an act of \emph{writing}. This inverts the
conventional view: the compiler is not the audience---the reader is.

We extend this insight: if programming is writing, then the
\emph{architecture} of a neural network is its \emph{grammar}. The choice
of attention mechanism, normalisation strategy, and feed-forward structure
are not merely engineering decisions---they are \emph{syntactic choices}
that constrain what the model can express, just as the choice between C
and Malbolge constrains what a programmer can communicate.

\subsection{Jungian Archetypes as Structural Priors}
\label{sec:jung}

Carl Gustav Jung proposed that beneath individual consciousness lies a
\emph{collective unconscious}---a shared psychic substrate populated by
\emph{archetypes}: recurring patterns of image, narrative, and affect
that manifest across cultures \citep{jung1968archetypes}. The archetypes
are not content but \emph{form}: the Hero is not any particular hero but
the \emph{pattern} by which heroism is recognised.

We operationalise 16~Jungian archetypes as geometric priors.
The selection requires justification: Jung himself never fixed a
canonical list. The \emph{Collected Works} treat archetypes as an
open set---``as many archetypes as there are typical situations in
life'' \citep[p.\,48]{jung1968archetypes}---while elaborating only a
handful in detail. Our 16 are drawn from three tiers of textual
authority:

\paragraph{Tier~1: Archetypes with dedicated treatment in Jung's
\emph{Collected Works}.}
Seven archetypes are extensively theorised by Jung himself:
\textbf{Self} (CW~9/2, \emph{Aion}),
\textbf{Shadow} (CW~9/2, ch.\,2--3),
\textbf{Anima} and \textbf{Animus} (CW~9/2, ch.\,3--4; CW~7),
\textbf{Persona} (CW~7, \emph{Two Essays on Analytical Psychology}),
\textbf{Wise Old Man}~/ \textbf{Spirit} (CW~9/1, ``The Phenomenology
of the Spirit in Fairy Tales''), and
\textbf{Divine Child} (CW~9/1, ``The Psychology of the Child
Archetype,'' joint with Ker\'enyi).
These are uncontroversial.

\paragraph{Tier~2: Archetypes named and characterised by Jung but
elaborated by his successors.}
Six archetypes are grounded in Jung's texts but received their full
theoretical articulation in the post-Jungian literature:
\textbf{Great Mother} and \textbf{Great Father} (named in CW~5,
\emph{Symbols of Transformation}; systematised by
Neumann~\citeyearpar{neumann1955great}),
\textbf{Hero} (present throughout CW~5; codified by
Campbell~\citeyearpar{campbell1949hero}),
\textbf{Trickster} (CW~9/1, ``On the Psychology of the Trickster
Figure''),
\textbf{Kore} (CW~9/1, ``The Psychological Aspects of the Kore''),
and \textbf{Puer Aeternus} (CW~5; given its canonical treatment by
von Franz).

\paragraph{Tier~3: Structural constructs used by Jung as organising
principles.}
Two entries are not archetypes in the strict imagistic sense but
structural categories that Jung uses to describe the psyche's
self-organisation:
\textbf{Senex} (the polar opposite of the Puer, named by Jung in
CW~5 as Saturn/Kronos and formalised by Hillman) and
\textbf{Quaternity} (CW~11, \emph{Psychology and Religion}; CW~12,
\emph{Psychology and Alchemy}---Jung's term for the fourfold structure
of psychic wholeness).

\paragraph{Why 16?}
The number is neither sacred nor arbitrary. It is the smallest set that
covers: (i)~the core ego--Self axis (Self, Shadow, Persona);
(ii)~the contrasexual pair (Anima, Animus); (iii)~the parental
imagos (Great Mother, Great Father); (iv)~the developmental spectrum
(Divine Child, Puer, Hero, Senex, Wise Old Man); (v)~the liminal and
structural figures (Trickster, Kore, Spirit, Quaternity). Any fewer
omits a structurally necessary pole; any more would require archetypes
that Jung himself did not identify, or would split existing ones into
sub-variants without clear theoretical warrant.

\paragraph{Why not more---or different?}
It should be noted that nothing in the architecture limits the number or
provenance of the curvature tensors. The metric $M = I + \sum_k w_k C_k$
accepts any number of tensors $C_k$ built from any source.
One could add Campbellian monomyth stages, Greimasian actantial roles,
Proppian narrative functions, Platonic forms, Aristotelian categories, or
curvature tensors derived from entirely non-Western ontologies such as
the \emph{gu\d{n}as} of S\=a\.{n}khya philosophy or the five
\emph{xing} of Chinese cosmology. Each would
produce a different deformation of the attention space---a different
grammar, in the sense of Section~\ref{sec:literate}.

The \textsc{DisattentionFormer} is therefore better understood not as a
\emph{Jungian} model but as a \emph{framework for embedding any
philosophical theory of form} as geometric structure in a Transformer.
The 16~Jungian archetypes are the instance we study; the architecture is
the general case. This generality is itself part of the argument: the
model does not \emph{understand} Jung any more than it would understand
Aristotle. It applies the covariance fingerprint of whatever passages
it is given. The philosophical content is in the researcher's choice of
tensors, not in the model's comprehension of them.

The 16 are enumerated below:

\begin{enumerate}[noitemsep]
  \item \textbf{Self} --- psychic totality, the mandala
  \item \textbf{Shadow} --- the unintegrated, the double
  \item \textbf{Anima} --- the feminine in the masculine psyche
  \item \textbf{Animus} --- the masculine in the feminine psyche
  \item \textbf{Great Mother} --- primordial feminine (nurturance/devouring)
  \item \textbf{Great Father} --- primordial masculine (law/tyranny)
  \item \textbf{Hero} --- ego confronting the unconscious
  \item \textbf{Wise Old Man} --- spirit as knowledge
  \item \textbf{Persona} --- the social mask
  \item \textbf{Puer Aeternus} --- the eternal youth
  \item \textbf{Senex} --- the old order, Saturn
  \item \textbf{Trickster} --- the boundary transgressor
  \item \textbf{Kore/Maiden} --- innocence and forced descent
  \item \textbf{Divine Child} --- pure potentiality
  \item \textbf{Spirit} --- orientation and illumination
  \item \textbf{Quaternity} --- the four functions, completeness
\end{enumerate}

Each archetype is converted into a curvature tensor $C_k \in
\mathbb{R}^{d \times d}$ via the following process: 10~representative
passages are encoded using a sentence transformer, their covariance
matrix is computed, projected to the model dimension via
Johnson--Lindenstrauss random projection, and trace-normalised
(\Cref{sec:tensors}).

\subsection{The Corpus as Architecture}
\label{sec:corpus}

The training corpus is \emph{not shuffled}. The sequence of 19~texts
follows a deliberate mythic progression:

\begin{enumerate}[noitemsep]
  \item \textbf{Creation phase:} Genesis $\to$ Rigveda Hymns $\to$ Popol Vuh
  \item \textbf{Heroic phase:} Iliad $\to$ Odyssey $\to$ Gilgamesh
  \item \textbf{Descent phase:} Homeric Hymns $\to$ Dante's Commedia
  \item \textbf{Transformation:} Grimm Tales $\to$ Blake (Songs, Prophetic, Marriage)
  \item \textbf{Individuation:} Vico's \textit{Scienza Nuova} $\to$ Joyce (Chamber Music $\to$ Dubliners $\to$ Portrait $\to$ Exiles $\to$ Ulysses)
\end{enumerate}

This ordering follows Vico's \textit{corsi e ricorsi}: theocratic $\to$
heroic $\to$ human $\to$ ricorso. The archetypal ``contamination'' of one
text over the next is the central mechanism---not a bug but the
architecture itself.

% ======================================================================
\section{The DisattentionFormer Architecture}
\label{sec:architecture}

\subsection{Overview}

The \textsc{DisattentionFormer} is a pre-norm Transformer
($d_{\text{model}} = 256$, $n_{\text{heads}} = 8$, $n_{\text{layers}} = 8$,
$d_{\text{ff}} = 1024$) augmented with four novel components:

\begin{enumerate}[noitemsep]
  \item \textbf{ArchetypalProjection} --- metric curvature via archetype
    tensors
  \item \textbf{SymbolAttention} --- self-attention in curved space
  \item \textbf{TensionFFN} --- opposing-pole feed-forward network
  \item \textbf{IndividuationNorm} --- depth-decaying normalisation
  \item \textbf{TranscendentFunction} --- archetype-gated synthesis
\end{enumerate}

The complete forward pass is:
\begin{align}
  \mathbf{x} &= \text{TokenEmb}(\text{ids}) + \text{PosEmb} \\
  M &= I + \sum_{k=1}^{16} w_k \, C_k, \quad
       w = \text{softmax}\!\left(\frac{\bar{\mathbf{x}} \cdot Q}{\sqrt{d}}\right) \\
  \mathbf{x}_{\text{curved}} &= \mathbf{x} \, M \\
  \text{(repeat } &8 \text{ blocks)} \notag \\
  \mathbf{x} &= \text{TranscendentFunction}(\mathbf{x}, w) \\
  \text{logits} &= \mathbf{x} \, W_{\text{emb}}^{\top} \quad \text{(weight tying)}
\end{align}

\subsection{Archetypal Curvature Tensors}
\label{sec:tensors}

For each archetype $k$, we select 10~canonical passages (e.g., for the
Hero: the call to adventure, the belly of the whale, the apotheosis).
These are encoded into embeddings $\mathbf{e}_i \in \mathbb{R}^{384}$ using
\texttt{all-MiniLM-L6-v2}, centered, and their covariance computed:
\begin{equation}
  C_k^{\text{raw}} = \frac{1}{9} \sum_{i=1}^{10}
    (\mathbf{e}_i - \bar{\mathbf{e}})(\mathbf{e}_i - \bar{\mathbf{e}})^{\top}
\end{equation}
This is projected from $\mathbb{R}^{384 \times 384}$ to
$\mathbb{R}^{256 \times 256}$ via a Johnson--Lindenstrauss random matrix
$R \in \mathbb{R}^{384 \times 256}$ (fixed seed), then trace-normalised:
\begin{equation}
  C_k = \frac{R^{\top} C_k^{\text{raw}} R}
             {\text{tr}(R^{\top} C_k^{\text{raw}} R) + \epsilon}
\end{equation}

The tensors are \textbf{not trainable}---they are registered as buffers.
This is deliberate: the collective unconscious precedes individual
learning.

\subsection{ArchetypalProjection}

The metric curvature $M$ warps the embedding space according to the
batch's archetypal composition:
\begin{equation}
  M = I + \sum_{k=1}^{16} w_k \, C_k, \quad
  \mathbf{x}_{\text{curved}} = \mathbf{x} \, M
\end{equation}
where $w_k = \text{softmax}(\bar{\mathbf{x}} \cdot \mathbf{q}_k / \sqrt{d})$
and $\mathbf{q}_k$ are learnable query vectors. The identity ensures that
the model starts as a standard Transformer and learns to \emph{deviate}
from flatness.

\subsection{SymbolAttention}

Standard self-attention computes scores as $QK^{\top}/\sqrt{d_h}$. We
insert the metric:
\begin{equation}
  \text{Scores} = \frac{(Q \, M) \, K^{\top}}{\sqrt{d_h}}
\end{equation}
This means that archetypally distant tokens can become close (or
closer) if the curvature tensor brings them together---a mythic token
in the Iliad might attend strongly to a Joycean token in \emph{Ulysses}
if both activate the Hero archetype.

\subsection{TensionFFN}

Each feed-forward layer computes two opposing views:
\begin{align}
  \mathbf{p} &= \text{GELU}(W_1^{+} \, \mathbf{x}) \quad \text{(thesis)} \\
  \mathbf{n} &= \text{GELU}(W_1^{-} \, (-\mathbf{x})) \quad \text{(antithesis)} \\
  \mathbf{y} &= W_2 \, [\mathbf{p} \,\|\, \mathbf{n}] \quad \text{(synthesis)}
\end{align}
The tension loss encourages the poles to diverge:
\begin{equation}
  \mathcal{L}_{\text{tension}} =
    1 - \cos(\mathbf{p}, \mathbf{n})
\end{equation}
maximised via negation in the total loss.

\subsection{IndividuationNorm}

Standard LayerNorm treats all layers identically. IndividuationNorm
interpolates between collective (normalised) and individual (raw)
representations based on depth:
\begin{equation}
  \text{IndivNorm}_i(\mathbf{x}) =
    (1 - r_i) \, \text{LayerNorm}(\mathbf{x}) + r_i \, \mathbf{x}, \quad
    r_i = \frac{i}{N - 1}
\end{equation}

Early layers ($r \approx 0$): strong normalisation, collective processing.
Deep layers ($r \approx 1$): minimal normalisation, individual expression.

\subsection{TranscendentFunction}

The final synthesis gate is modulated by the active archetypal field:
\begin{align}
  \mathbf{s} &= \tanh(W_s \, \mathbf{x}) \\
  g &= \sum_{k=1}^{16} w_k \, \sigma(W_{g,k} \, \mathbf{x}) \\
  \text{out} &= g \cdot \mathbf{s} + (1 - g) \cdot \mathbf{x}
\end{align}

% ======================================================================
\section{Training}
\label{sec:training}

\subsection{Loss Function}

The total loss is:
\begin{equation}
  \mathcal{L} = \mathcal{L}_{\text{ce}}
    + \alpha \, \mathcal{L}_{\text{constellation}}
    - \beta \, \mathcal{L}_{\text{tension}}
    - \gamma \, \mathcal{L}_{\text{individuation}}
\end{equation}
with $\alpha = 0.10$, $\beta = 0.05$, $\gamma = 0.02$.

\begin{itemize}[noitemsep]
  \item $\mathcal{L}_{\text{ce}}$: standard cross-entropy (next-token prediction)
  \item $\mathcal{L}_{\text{constellation}}$: tokens with similar archetypal
    weights should have similar representations (cosine alignment)
  \item $\mathcal{L}_{\text{tension}}$: maximise divergence between TensionFFN
    poles (the model must maintain internal conflict)
  \item $\mathcal{L}_{\text{individuation}}$: maximise representation divergence
    from the batch centroid (each token should become itself)
\end{itemize}

\subsection{Training Configuration}

Both models were trained on an Apple~M3 MacBook Pro (18\,GB) using MPS
backend in float32 (no AMP). Identical hyperparameters except architecture:

\begin{table}[H]
\centering
\caption{Training configuration (shared).}
\label{tab:training-config}
\begin{tabular}{ll}
\toprule
\textbf{Hyperparameter} & \textbf{Value} \\
\midrule
Optimiser       & AdamW ($\beta_1 = 0.9$, $\beta_2 = 0.95$) \\
Learning rate   & $1 \times 10^{-4}$ \\
Schedule        & Linear warmup (2000 steps) + cosine decay \\
Weight decay    & 0.1 (dim $\geq 2$), 0.0 (biases, norms) \\
Batch size      & 4 \\
Sequence length & 512 tokens \\
Gradient clip   & 1.0 (max norm) \\
Epochs          & 10 \\
Precision       & float32 \\
\bottomrule
\end{tabular}
\end{table}

\subsection{Corpus Details}

The tokeniser is word-level with 62\,417~types built from the corpus.
The 19~source texts total approximately~X million words (see
\Cref{sec:corpus}).

% ======================================================================
\section{Baseline: Parameter-Matched Transformer}
\label{sec:baseline}

To isolate the contribution of the DisattentionFormer's novel components,
we train a standard pre-norm Transformer with matched parameter count:

\begin{table}[H]
\centering
\caption{Architectural comparison.}
\label{tab:arch-comparison}
\begin{tabular}{lcc}
\toprule
\textbf{Component} & \textbf{DisattentionFormer} & \textbf{Baseline} \\
\midrule
$d_{\text{model}}$          & 256       & 272       \\
$n_{\text{heads}}$          & 8         & 8         \\
$d_{\text{head}}$           & 32        & 34        \\
$n_{\text{layers}}$         & 8         & 12        \\
$d_{\text{ff}}$             & 1024      & 1088      \\
max seq.\ len.              & 512       & 512       \\
Dropout                     & 0.1       & 0.1       \\
\midrule
Archetype Tensors ($C_k$)   & \checkmark & ---      \\
SymbolAttention              & \checkmark & ---      \\
TensionFFN                   & \checkmark & ---      \\
IndividuationNorm            & \checkmark & ---      \\
TranscendentFunction         & \checkmark & ---      \\
Auxiliary Losses             & \checkmark & ---      \\
\midrule
Total parameters            & $\approx$27,722,496 & $\approx$27,783,712 \\
\bottomrule
\end{tabular}
\end{table}

The baseline uses standard multi-head self-attention, GELU feed-forward,
pre-norm LayerNorm, and weight tying. It is trained with identical
hyperparameters, the same corpus, and the same tokeniser. The only
loss is cross-entropy.

Note that the baseline has more layers (12 vs.\ 8) and a slightly wider
hidden dimension (272 vs.\ 256) to match the parameter count. This
\emph{favours} the baseline: more depth typically improves language
modelling \citep{devlin2019bert}, and the baseline allocates its entire
parameter budget to language modelling rather than splitting it with
auxiliary components.

% ======================================================================
\section{Results}
\label{sec:results}

% These values will be filled in after running benchmark.py
% Use \begin{table}[h]
\centering
\caption{Quantitative comparison on the mythic--Joycean corpus (3 models).
All models trained for 10 epochs with identical optimisation schedule.}
\label{tab:benchmark}
\begin{tabular}{lccc}
\toprule
\textbf{Metric} & \textbf{DisattentionFormer V1} & \textbf{DisattentionFormer V2 (Titans)} & \textbf{Baseline Transformer} \\
\midrule
Parameters & 27,722,496 & 25,690,112 & 27,783,984 \\
$d_{{\text{{model}}}}$ & 256 & 256 & 272 \\
Layers & 8 & 6 & 12 \\
Heads & 8 & 8 & 8 \\
$d_{{\text{{ff}}}}$ & 1024 & 1024 & 1088 \\
\midrule
Perplexity & 845.84 & 522.27 & 339.01 \\
CE Loss & 6.7403 & 6.2582 & 5.8260 \\
\midrule
PPL (ctx=128) & 628.53 & 432.84 & 365.07 \\
PPL (ctx=256) & 658.58 & 457.91 & 362.77 \\
PPL (ctx=512) & 680.57 & 466.09 & 367.8 \\
\midrule
Type-Token Ratio & 0.1128 & 0.1526 & 0.1733 \\
Repetition Rate & 0.2174 & 0.2066 & 0.1924 \\
Mean Entropy & 6.1664 & 5.9166 & 5.9645 \\
\midrule
Tokens/s & 25878 & 2960 & 19824 \\
\midrule
Archetype Tensors & $\checkmark$ & $\checkmark$ & --- \\
Neural Memory & --- & $\checkmark$ & --- \\
\bottomrule
\end{tabular}
\end{table}
 to include
% the auto-generated table.

\subsection{Perplexity}

\begin{table}[H]
\centering
\caption{Validation set metrics after 10 epochs of training.}
\label{tab:results}
\begin{tabular}{lcc}
\toprule
\textbf{Metric} & \textbf{DisattentionFormer} & \textbf{Baseline} \\
\midrule
Val.\ CE Loss       & \texttt{TBD} & \texttt{TBD} \\
Val.\ Perplexity    & \texttt{TBD} & \texttt{TBD} \\
Mean Token Entropy   & \texttt{TBD} & \texttt{TBD} \\
\bottomrule
\end{tabular}
\end{table}

\subsection{Training Dynamics}

The DisattentionFormer's training exhibited several distinctive patterns:

\begin{itemize}[noitemsep]
  \item \textbf{Loss decomposition:} Final epoch averages---$\mathcal{L}_{\text{ce}} \approx 6.5$,
    $\mathcal{L}_{\text{constellation}} \approx -0.025$,
    $\mathcal{L}_{\text{tension}} \approx 0.97$,
    $\mathcal{L}_{\text{individuation}} = 10.0$ (clamped).
  \item \textbf{Individuation saturation:} $\mathcal{L}_{\text{individuation}}$
    immediately saturated at its clamp value of 10.0, suggesting the model
    maximises representation divergence to the limit of what we allow.
  \item \textbf{Tension convergence:} $\mathcal{L}_{\text{tension}}$ increased
    monotonically from $\approx$0.88 to $\approx$0.97, indicating the
    TensionFFN poles are learning to diverge throughout training.
  \item \textbf{Dominant archetype shift:} Great Mother dominated early training
    (creation texts), shifting to Trickster by the middle epochs (heroic and
    Joycean texts).
\end{itemize}

\begin{figure}[H]
\centering
% Placeholder for training curves
% \input{figures/training_curves.pgf}
\fbox{\parbox{0.9\textwidth}{\centering\vspace{3em}
  Training loss curves: DisattentionFormer (total, CE, auxiliary)
  vs.\ Baseline (CE only). To be generated from training logs.
  \vspace{3em}}}
\caption{Training loss over 10 epochs (7280 steps per model).}
\label{fig:training-curves}
\end{figure}

\subsection{Generation Quality}

\begin{table*}[t]
\centering
\caption{Generation samples from thematic prompts (temperature $= 0.8$,
  nucleus $p = 0.9$, 100 tokens).}
\label{tab:generation}
\small
\begin{tabular}{p{2.5cm}p{5.5cm}p{5.5cm}}
\toprule
\textbf{Prompt} & \textbf{DisattentionFormer} & \textbf{Baseline} \\
\midrule
``In the beginning'' & \texttt{TBD} & \texttt{TBD} \\
\midrule
``Stately, plump'' & \texttt{TBD} & \texttt{TBD} \\
\midrule
``I will not serve'' & \texttt{TBD} & \texttt{TBD} \\
\bottomrule
\end{tabular}
\end{table*}

\subsection{Archetype Activation Analysis}

% DisattentionFormer-specific analysis

\begin{table}[H]
\centering
\caption{Archetype activation frequencies across thematic prompts
  (DisattentionFormer). Each cell shows the fraction of generated
  tokens for which the given archetype was dominant.}
\label{tab:archetypes}
\begin{tabular}{lc}
\toprule
\textbf{Archetype} & \textbf{Frequency} \\
\midrule
Self           & \texttt{TBD} \\
Shadow         & \texttt{TBD} \\
Anima          & \texttt{TBD} \\
Great Mother   & \texttt{TBD} \\
Great Father   & \texttt{TBD} \\
Hero           & \texttt{TBD} \\
Trickster      & \texttt{TBD} \\
\emph{(others)}& \texttt{TBD} \\
\bottomrule
\end{tabular}
\end{table}

% ======================================================================
\section{Discussion}
\label{sec:discussion}

\subsection{Architecture as Grammar}

The central thesis of this paper is that the \emph{architecture} of a
neural network is a \emph{language}---not metaphorically, but structurally.
Just as the C compiler enforces a particular relationship between the
programmer's intention and the machine's execution, the Transformer's
attention mechanism enforces a particular relationship between tokens.

The DisattentionFormer makes this explicit. Its components are not
arbitrary engineering choices but deliberate \emph{syntactic} commitments:

\begin{itemize}[noitemsep]
  \item \textbf{ArchetypalProjection} is a \emph{declension}---it inflects
    the embedding space according to archetypal context, like how the
    genitive case in Latin inflects a noun according to its possessor.
  \item \textbf{TensionFFN} is a \emph{dialectic}---thesis ($+\mathbf{x}$)
    and antithesis ($-\mathbf{x}$) coexist without resolution, like the
    opposing clauses of a Joycean sentence.
  \item \textbf{IndividuationNorm} is a \emph{voice}---early layers speak
    in a collective register (normalised), while deep layers speak
    individually (raw).
  \item \textbf{TranscendentFunction} is a \emph{conjugation}---it modulates
    the output according to tense (which archetypes are active), producing
    the appropriate inflection of meaning.
\end{itemize}

\subsection{A Compiler for the Collective Unconscious}
\label{sec:compiler}

The most revealing result of this work is not quantitative but structural.
We built what might be called a \emph{compiler for the collective
unconscious}: an architecture that takes Jungian theory as source code,
covariance geometry as intermediate representation, and curved attention
as executable output. And it compiled. It took our rich symbolic input and
reduced it to covariance matrices. It took mythic progression and reduced
it to training-order effects. It took the tension between thesis and
antithesis and reduced it to cosine distance.

Consider the analogy at three levels of the compilatory chain:

\begin{enumerate}
  \item \textbf{The C compiler} translates human intent to machine
    execution. It does not \emph{understand} \texttt{"Hello world"};
    it parses syntax, resolves symbols, emits instructions. But the
    structure of C---types, functions, return codes---shapes \emph{what
    can be said}. The grammar is not understanding, but it constrains
    communication.
  \item \textbf{The vanilla Transformer} translates statistical
    regularities into predictions. It does not \emph{understand} Joyce;
    it parses co-occurrence patterns across 62\,417 tokens. The
    architecture---attention, FFN, LayerNorm---shapes \emph{what can be
    learned}. The architecture is not understanding, but it constrains
    generation.
  \item \textbf{The DisattentionFormer} encodes philosophical structure
    into the architecture itself. Jungian archetypes as curvature,
    opposing poles in the FFN, depth-decaying norms. We give it every
    advantage: the theory is explicit, the geometry is precise, the
    corpus is ordered. And what happens? The model still operates at
    the level of lexical fingerprints. It activates ``Great Father''
    for Genesis not because it grasps patriarchal divinity but because
    the words overlap. The architecture is not understanding---it is
    a more elaborate grammar.
\end{enumerate}

At each level, the same operation occurs: \textbf{faithful translation
without comprehension}. The IOCCC connection from Section~\ref{sec:introduction}
is not decorative but structural. Obfuscated C code and LLM output
share the same property: \emph{surface complexity emerging from mechanical
simplicity}, with the human observer supplying the meaning. The circle
of underscores that computes $\pi$ is, underneath, just \texttt{--F} and
\texttt{OO++}. The Transformer that generates Joycean prose is, underneath,
just matrix multiplications and softmax distributions.

This is not a limitation to lament but a fact to understand.
The value of understanding it---\emph{really} understanding it, as a
programmer understands what \texttt{int main()} does---is greater than
the value of prompting a black box to produce plausible text.

\subsection{What the Tensors Cannot See}
\label{sec:tensors-cannot}

The archetypal tensors $C_k$ are covariance matrices of sentence
embeddings. They capture the \emph{shape} (directions of variance) of
how each archetype's seed passages spread in embedding space. This is
a statistical fingerprint, not a conceptual representation.

Concretely, the model cannot recognise that:

\begin{itemize}[noitemsep]
  \item \textbf{Shadow and Persona are dialectically related.} They are
    stored as two independent curvature matrices. There is no mechanism
    to represent that the Shadow is what the Persona conceals---that
    they are the same structure seen from opposite sides.
  \item \textbf{God in Genesis is a Great Father \emph{figure}.} The
    model activates the Great Father tensor for Genesis because the
    seed passages include ``I am the Lord your God'' and the words
    overlap. If a passage described a tyrannical father using
    entirely novel vocabulary, the tensor would not activate.
  \item \textbf{The Hero is a narrative pattern, not a word cluster.}
    The Hero archetype is defined by its \emph{function}---descent,
    confrontation, return---which requires tracking narrative dynamics
    across a sequence. The model computes a sequence mean and takes
    a dot product. It sees hero-\emph{words}, not heroism.
\end{itemize}

The learnable query vectors $\mathbf{q}_k$ offer no escape. They learn
statistical correlations (``when I activate this curvature, training loss
decreases''), but they cannot learn \emph{why}.

\subsection{Limitations}

\begin{itemize}[noitemsep]
  \item \textbf{Archetype tensor rank:} With only 10~passages per archetype,
    the covariance matrices are rank-deficient ($\text{rank} \leq 9$ in
    $\mathbb{R}^{256 \times 256}$). The tensors capture surface
    statistics, not the deep functional role of archetypes.
  \item \textbf{Weight uniformity:} Archetype weights remain nearly uniform
    ($\approx 0.0625 = 1/16$) throughout training, suggesting the model
    has not learned strong archetypal differentiation.
  \item \textbf{Individuation clamp:} $\mathcal{L}_{\text{individuation}}$
    saturates at the clamp value immediately, indicating the loss scale
    needs retuning.
  \item \textbf{No inter-archetype structure:} The 16 tensors are linearly
    combined with no relational modelling. Jungian theory posits
    compensatory relationships (Puer--Senex, Shadow--Persona,
    Anima--ego)---none of which can be captured by independent
    weighting.
  \item \textbf{Single-machine training:} All experiments were conducted on
    a single Apple~M3 with 18\,GB RAM, limiting batch size and precision.
\end{itemize}

\subsection{On the Originality of Philosophical Tensors}

A literature search of arXiv and related repositories revealed no prior
work that constructs text-derived covariance tensors from philosophical or
literary passages and uses them as fixed geometric priors to deform
attention space. Related work in hyperbolic embeddings
\citep{nickel2017poincare} and geometric deep learning
\citep{bronstein2017geometric} operates in continuous curvature spaces but
does not derive curvature from \emph{textual content}. The specific
combination of Jungian theory, covariance geometry, and Transformer
architecture appears to be novel.

% ======================================================================
\section{Conclusion}
\label{sec:conclusion}

We have demonstrated three things:

\begin{enumerate}
  \item That a Transformer \emph{can} be designed as a structured grammar
    encoding psychological and mythic theory---the
    \textsc{DisattentionFormer} is an existence proof.
  \item That such an architecture, despite its rich philosophical
    scaffolding, still operates by the same mechanism as every other
    model in the compilatory chain: faithful translation without
    comprehension. It parses lexical co-occurrence where we intended
    archetypal recognition. It computes cosine distances where we
    intended dialectical tension.
  \item That \emph{understanding this} is the point. The
    \textsc{DisattentionFormer} does not outperform a vanilla Transformer
    on perplexity---nor does it need to. Its value lies in making
    visible what is normally hidden: that the architecture is the
    message, and that every model, no matter how elaborate, is
    ultimately a parser.
\end{enumerate}

Knuth argued that programs should be written for human readers, with
machine execution secondary. We extend this: neural architectures should
be \emph{read}, not merely prompted. The difference between a programmer
who understands \texttt{int main()} and a user who types
\texttt{"write me a poem"} is the difference between literacy and
consumption. Large language models are not magic. They are compilers
without a spec sheet---and the path forward is to write the spec sheet,
not to marvel at the output.

The act of developing software is an act of translation. The act of
designing a neural architecture is an act of grammar. And the act of
training on an unshuffled mythic corpus---from Genesis to
\textit{Ulysses}---is an act of reading, performed by a machine that
cannot read.

\medskip
\noindent
\textit{Shut your eyes and see.}

\hfill --- James Joyce, \textit{Ulysses} (3.9)

% ======================================================================
% References
% ======================================================================
\bibliography{references}

\end{document}
